\documentclass[10pt]{article}

\usepackage{amsmath, amssymb, amsthm}
\usepackage{geometry}
\usepackage{multicol}
\usepackage{enumitem}
\usepackage{xcolor}
\usepackage[most]{tcolorbox}
\usepackage{bm}
\usepackage{titlesec}

\geometry{margin=0.52in}
\setlength{\columnsep}{0.22in}
\setlength{\parindent}{0pt}
\setlength{\parskip}{2pt}

\titleformat{\section}
{\large\bfseries\color{blue}}
{\thesection}{0.4em}{}

\titleformat{\subsection}
{\bfseries\color{teal!70!black}}
{\thesubsection}{0.4em}{}

% ---------- Color Boxes ----------

\newtcolorbox{definitionbox}{
colback=blue!5,
colframe=blue!60!black,
boxrule=0.8pt,
arc=4pt,
left=4pt,right=4pt,top=4pt,bottom=4pt
}

\newtcolorbox{theorembox}{
colback=green!5,
colframe=green!50!black,
boxrule=0.8pt,
arc=4pt,
left=4pt,right=4pt,top=4pt,bottom=4pt
}

\newtcolorbox{formulabox}{
colback=yellow!10,
colframe=orange!80!black,
boxrule=0.8pt,
arc=4pt,
left=4pt,right=4pt,top=4pt,bottom=4pt
}

\title{\vspace{-1cm}STAT 545 Final Exam Notesheet}
\date{}

\begin{document}

\maketitle
\vspace{-1.1cm}

\begin{multicols}{2}

%------------------------------------------------
\section{Stochastic Processes}

\begin{definitionbox}
\textbf{Stochastic Process}

A stochastic process is a collection of random variables
\[
\{X_t : t \in I\}
\]
where \(I\) is the index set (usually time) and \(S\) is the state space.
\end{definitionbox}

%------------------------------------------------
\section{Markov Chains}

\begin{definitionbox}
\textbf{Markov Property}
\[
P(X_{n+1}=j \mid X_0,\dots,X_n=i)
=
P(X_{n+1}=j \mid X_n=i)
\]
\end{definitionbox}

\subsection{Transition Matrix}

\begin{formulabox}
\[
P_{ij} = P(X_{n+1}=j \mid X_n=i)
\]
\[
P =
\begin{pmatrix}
p_{11} & p_{12} & \dots \\
p_{21} & p_{22} & \dots \\
\vdots & \vdots & \ddots
\end{pmatrix}
\]
\end{formulabox}

Properties:
\begin{itemize}[leftmargin=12pt,itemsep=1pt]
\item \(p_{ij}\ge 0\)
\item \(\sum_j p_{ij}=1\)
\end{itemize}

\subsection{Time Homogeneous}

\[
P(X_{n+1}=j \mid X_n=i)=P(X_1=j \mid X_0=i)
\]

%------------------------------------------------
\section{\(n\)-Step Transition Probabilities}

\begin{formulabox}
\[
P_{ij}^{(n)} = P(X_n=j \mid X_0=i)
\]
\[
P^{(n)} = P^n
\]
\end{formulabox}

\begin{theorembox}
\textbf{Chapman--Kolmogorov}
\[
P_{ij}^{(m+n)}
=
\sum_k P_{ik}^{(m)}P_{kj}^{(n)}
\]
\end{theorembox}

%------------------------------------------------
\section{Distribution Evolution}

\begin{formulabox}
If initial distribution is
\[
\alpha = (\alpha_1,\dots,\alpha_k)
\]
then
\[
P(X_n=\cdot)=\alpha P^n
\]
\end{formulabox}

%------------------------------------------------
\section{Stationary Distributions}

\begin{definitionbox}
\[
\pi=\pi P
\]
\[
\pi_j=\sum_i \pi_i P_{ij}
\]
with
\[
\sum_i \pi_i=1
\]
\end{definitionbox}

%------------------------------------------------
\section{Regular / Ergodic Chains}

\begin{definitionbox}
A matrix \(P\) is \textbf{regular} if
\[
P^n>0
\]
for some \(n\).
\end{definitionbox}

\begin{theorembox}
If \(P\) is regular:
\begin{itemize}[leftmargin=12pt,itemsep=1pt]
\item unique stationary distribution
\item limiting distribution exists
\item rows of \(P^n\) converge to \(\pi\)
\end{itemize}
\end{theorembox}

\begin{theorembox}
Finite chain is \textbf{ergodic} if
\begin{itemize}[leftmargin=12pt,itemsep=1pt]
\item irreducible
\item aperiodic
\end{itemize}
Then
\[
P_{ij}^{(n)}\to \pi_j
\]
\end{theorembox}

%------------------------------------------------
\section{Communication, Recurrence, Period}

\begin{definitionbox}
State \(j\) is accessible from \(i\) if
\[
P_{ij}^{(n)}>0
\]
for some \(n\ge 1\).
\end{definitionbox}

Communicating: accessible both ways.

Irreducible: all states communicate.

\begin{definitionbox}
State \(j\) is recurrent if
\[
P(\text{return to }j)=1
\]
Equivalent:
\[
\sum_{n=0}^\infty P_{jj}^{(n)}=\infty
\]
Transient if
\[
\sum_{n=0}^\infty P_{jj}^{(n)}<\infty
\]
\end{definitionbox}

\begin{formulabox}
\[
d(i)=\gcd\{n>0 : P_{ii}^{(n)}>0\}
\]
Aperiodic if \(d(i)=1\).
\end{formulabox}

%------------------------------------------------
\section{Detailed Balance / Reversibility}

\begin{formulabox}
\[
\pi(i)P(i,j)=\pi(j)P(j,i)
\]
\end{formulabox}

Detailed balance \(\Rightarrow\) \(\pi\) stationary.

Reversible chain \(\iff\) detailed balance holds for some \(\pi\).

%------------------------------------------------
\section{Birth--Death Chains}

\[
P(i,i+1)=p_i,\qquad P(i,i-1)=q_i
\]

\begin{formulabox}
\[
\frac{\pi(i+1)}{\pi(i)}=\frac{p_i}{q_{i+1}}
\]
\end{formulabox}

%------------------------------------------------
\section{Reflecting Random Walks}

Interior transitions:
\[
P(i,i+1)=p,\qquad P(i,i-1)=1-p
\]

Typical reflecting boundaries:
\[
P(0,1)=1,\qquad P(n,n-1)=1
\]

\begin{formulabox}
Let
\[
r=\frac{p}{1-p}
\]
Then from detailed balance,
\[
\pi(i+1)=r\,\pi(i)
\]
up to boundary adjustments / normalization.
\end{formulabox}

For finite path with no self-loops: usually irreducible, period \(2\).

%------------------------------------------------
\section{Mean Return Time}

\begin{formulabox}
For positive recurrent state \(i\),
\[
E[T_i]=\frac{1}{\pi(i)}
\]
\end{formulabox}

Long-run fraction of time in state \(i\) is \(\pi(i)\).

%------------------------------------------------
\section{Hitting Times / Hitting Probabilities}

\begin{definitionbox}
For set \(A\),
\[
\tau_A=\inf\{t\ge 0: X_t\in A\}
\]
For a single state \(j\),
\[
\tau_j=\inf\{t\ge 0: X_t=j\}
\]
\end{definitionbox}

\begin{formulabox}
Hitting probability:
\[
h_i=P_i(\tau_A<\infty)
\]
satisfies
\[
h_i=\sum_j P_{ij}h_j
\]
for \(i\notin A\), with boundary
\[
h_i=1,\quad i\in A
\]
\end{formulabox}

\begin{formulabox}
Expected hitting time:
\[
m_i=E_i[\tau_A]
\]
satisfies
\[
m_i=1+\sum_j P_{ij}m_j
\]
for \(i\notin A\), with
\[
m_i=0,\quad i\in A
\]
\end{formulabox}

Try exponential form \(h_i=\lambda^i\) for nearest-neighbor walks.

%------------------------------------------------
\section{Recurrence Classification on \(\mathbb{N}\)}

For reflecting walk on \(\mathbb{N}\) with interior moves
\[
P(i,i+1)=p,\qquad P(i,i-1)=1-p,\qquad P(0,1)=1
\]

\begin{formulabox}
\[
r=\frac{p}{1-p}
\]
\[
p<\frac12 \Rightarrow \text{positive recurrent}
\]
\[
p=\frac12 \Rightarrow \text{null recurrent}
\]
\[
p>\frac12 \Rightarrow \text{transient}
\]
\end{formulabox}

Characteristic equation for return probs:
\[
p\lambda^2-\lambda+(1-p)=0
\]

%------------------------------------------------
\section{Two-State Markov Chains}

\[
P=
\begin{pmatrix}
1-p & p\\
q & 1-q
\end{pmatrix}
\qquad
\alpha=(\alpha_0,\alpha_1)
\]

\subsection{Two-Step Transition}

\begin{formulabox}
\[
P^2=
\begin{pmatrix}
(1-p)^2+pq & p(2-p-q)\\
q(2-p-q) & (1-q)^2+pq
\end{pmatrix}
\]
\end{formulabox}

\subsection{Expectation of the State}

Since \(X_n\in\{0,1\}\),
\[
E[X_n]=P(X_n=1)
\]
\[
E[X_2]=(\alpha P^2)_1
\]

\subsection{Closed Form for \(P^n\)}

\begin{formulabox}
\[
P^n=
\frac{1}{p+q}
\begin{pmatrix}
q+p(1-p-q)^n & p-p(1-p-q)^n\\
q-q(1-p-q)^n & p+q(1-p-q)^n
\end{pmatrix}
\]
\end{formulabox}

\[
\pi=
\left(\frac{q}{p+q},\frac{p}{p+q}\right)
\]

\[
\lim_{n\to\infty}P^n=
\begin{pmatrix}
\pi_0 & \pi_1\\
\pi_0 & \pi_1
\end{pmatrix}
\]

\begin{formulabox}
\[
E[X_n]
=
\frac{p}{p+q}
+
\left(\alpha_1-\frac{p}{p+q}\right)(1-p-q)^n
\]
\end{formulabox}

%------------------------------------------------
\section{Random Walks on Graphs}

\begin{definitionbox}
A random walk on a graph moves uniformly to a neighboring vertex.
\end{definitionbox}

\[
P_{ij}=
\frac{1}{\deg(i)}
\quad \text{if } i\sim j
\]

Key steps:
\begin{enumerate}[leftmargin=14pt,itemsep=1pt]
\item Compute degree of each vertex.
\item Normalize each row by degree.
\end{enumerate}

\begin{formulabox}
Stationary distribution for simple random walk on undirected graph:
\[
\pi(i)=\frac{\deg(i)}{\sum_j \deg(j)}
\]
\end{formulabox}

\subsection{Period in Graph Walks}

\begin{formulabox}
\[
d(i)=\gcd\{n\ge 1 : P_{ii}^{(n)}>0\}
\]
\end{formulabox}

Facts:
\begin{itemize}[leftmargin=12pt,itemsep=1pt]
\item bipartite graph \(\Rightarrow\) period \(2\)
\item self-loop often makes chain aperiodic
\end{itemize}

%------------------------------------------------
\section{Absorbing Markov Chains}

\begin{definitionbox}
A state is absorbing if
\[
P_{ii}=1
\]
\end{definitionbox}

Reorder states so
\[
P=
\begin{pmatrix}
Q & R\\
0 & S
\end{pmatrix}
\]
where \(Q\): transient \(\to\) transient, \(R\): transient \(\to\) closed/absorbing.

\subsection{Fundamental Matrix}

\begin{theorembox}
\[
F=(I-Q)^{-1}
\]
\end{theorembox}

Interpretation:
\[
F_{ij}=\text{expected visits to transient state }j\text{ starting from }i
\]

\subsection{Expected Time to Absorption}

\begin{formulabox}
\[
E[T\mid X_0=i]=\sum_j F_{ij}
\]
\end{formulabox}

\subsection{Absorption Probabilities}

\begin{formulabox}
\[
B=FR
\]
\end{formulabox}

\(B_{ij}\): probability of eventual absorption in closed/absorbing state \(j\).

\subsection{Limiting Matrix}

\begin{theorembox}
If closed classes are aperiodic,
\[
\lim_{n\to\infty}P^n=
\begin{pmatrix}
0 & (I-Q)^{-1}R\Pi\\
0 & \Pi
\end{pmatrix}
\]
\end{theorembox}

\(\Pi\): stationary distributions of closed classes.

%------------------------------------------------
\section{Estimating Transition Matrices}

Suppose transition counts \(N_{ij}\) are observed.

\subsection{MLE}

\begin{formulabox}
\[
\hat P_{ij}=
\frac{N_{ij}}{\sum_j N_{ij}}
\]
\end{formulabox}

(row-normalized counts)

\subsection{Laplace Smoothing}

\begin{formulabox}
\[
\hat p_i=
\frac{N_i+1}{\sum_j (N_j+1)}
\]
\end{formulabox}

Prevents zero probabilities.

\subsection{Dirichlet Smoothing}

Prior:
\[
p\sim Dir(\alpha\mu)
\]

Posterior mean:
\begin{formulabox}
\[
\tilde p_i=
\frac{N_i+\alpha\mu_i}{N+\alpha}
\]
\end{formulabox}

%------------------------------------------------
\section{Poisson Process}

\begin{definitionbox}
For homogeneous Poisson process \((N(t))\):
\begin{itemize}[leftmargin=12pt,itemsep=1pt]
\item \(N(0)=0\)
\item independent increments
\item stationary increments
\item \(N(t)\sim \text{Poisson}(\lambda t)\)
\end{itemize}
\end{definitionbox}

\begin{formulabox}
\[
P(N(t)=k)=e^{-\lambda t}\frac{(\lambda t)^k}{k!}
\]
\[
E[N(t)]=\lambda t,\qquad Var(N(t))=\lambda t
\]
\end{formulabox}

\subsection{MLE for Poisson Rate}

\begin{formulabox}
\[
\hat\lambda=\frac{N(T)}{T}
\]
\end{formulabox}

\subsection{Exponential Waiting Times}

\[
T_1\sim Exp(\lambda)
\]
\[
P(T_1>t)=e^{-\lambda t}
\]
\[
E[T_1]=\frac{1}{\lambda}
\]

If waiting for \(m\) additional arrivals,
\[
E[\text{wait time}]=\frac{m}{\lambda}
\]

\subsection{Thinning}

If each event is observed independently with probability \(q\),
\[
\lambda_{\text{obs}}=q\lambda,\qquad \lambda=\frac{\lambda_{\text{obs}}}{q}
\]

%------------------------------------------------
\section{Nonhomogeneous Poisson Process}

Intensity \(\lambda(t)\) varies with time.

\begin{formulabox}
\[
E[N(b)-N(a)]=\int_a^b \lambda(t)\,dt
\]
\[
N(b)-N(a)\sim \text{Poisson}\!\left(\int_a^b \lambda(t)\,dt\right)
\]
\end{formulabox}

Probability of no arrivals on \([0,t]\):
\[
P(N(t)=0)=\exp\left(-\int_0^t \lambda(s)\,ds\right)
\]

First arrival survival:
\[
P(T_1>t)=\exp\left(-\int_0^t \lambda(s)\,ds\right)
\]

%------------------------------------------------
\section{Brownian Motion}

\begin{definitionbox}
Standard Brownian motion \((B_t)\):
\begin{itemize}[leftmargin=12pt,itemsep=1pt]
\item \(B_0=0\)
\item independent increments
\item \(B_t-B_s\sim N(0,t-s)\) for \(t>s\)
\end{itemize}
\end{definitionbox}

\[
E[B_t]=0,\qquad Var(B_t)=t
\]
\[
Cov(B_s,B_t)=\min(s,t)
\]

%------------------------------------------------
\section{Brownian Motion with Drift}

\begin{formulabox}
\[
X_t=\mu t+\sigma B_t
\]
\[
E[X_t]=\mu t,\qquad Var(X_t)=\sigma^2 t
\]
\end{formulabox}

If observed at equal spacing \(\Delta t\),
\[
\hat\mu=\frac{1}{n\Delta t}\sum_{k=1}^n (X_{t_k}-X_{t_{k-1}})
\]

For 2D:
\[
Z_t=\mu t+\Sigma^{1/2}B_t,
\qquad
\hat\mu=\frac{Z_{t_n}-Z_{t_0}}{t_n-t_0}
\]

%------------------------------------------------
\section{Brownian Bridge}

\begin{definitionbox}
Standard Brownian bridge:
\[
W_t=B_t-tB_1,\qquad 0\le t\le 1
\]
Equivalent in distribution to \((B_t\mid B_1=0)\).
\end{definitionbox}

\[
E[W_t]=0
\]
\[
Cov(W_s,W_t)=\min(s,t)-st
\]

For \(X_t=\mu t+\sigma B_t\),
\[
(X_t\mid X_1=0)\overset{d}= \sigma(B_t-tB_1)
\]
with covariance
\[
\sigma^2(\min(s,t)-st)
\]

%------------------------------------------------
\section{Conditional Normal}

\begin{formulabox}
If \((X,Y)\) is jointly normal, then
\[
X\mid Y=y \sim
N\left(
\mu_x+\frac{\sigma_{xy}}{\sigma_y^2}(y-\mu_y),
\;
\sigma_x^2-\frac{\sigma_{xy}^2}{\sigma_y^2}
\right)
\]
\end{formulabox}

Useful with Brownian motion:
\[
(B_t,B_1)\sim N\!\left(
\begin{pmatrix}0\\0\end{pmatrix},
\begin{pmatrix}t&t\\ t&1\end{pmatrix}
\right)
\]

%------------------------------------------------
\section{Random Walk Drift in 2D}

For steps \(E,W,N,S\) with probabilities \(p_E,p_W,p_N,p_S\),

\begin{formulabox}
\[
E[\Delta Z]=(p_E-p_W,\; p_N-p_S)
\]
\end{formulabox}

If counts are \(N_E,N_W,N_N,N_S\) over \(n\) steps,
\[
\hat p_E=\frac{N_E}{n},\quad
\hat p_W=\frac{N_W}{n},\quad
\hat p_N=\frac{N_N}{n},\quad
\hat p_S=\frac{N_S}{n}
\]

\[
\hat\delta=
\left(
\frac{N_E-N_W}{n},
\frac{N_N-N_S}{n}
\right)
\]

Sample mean increment:
\[
\bar\Delta=\frac1n\sum_{k=1}^n (Z_k-Z_{k-1})
\]

If each step lasts \(\Delta t\),
\[
\hat\mu=\frac{\hat\delta}{\Delta t}
\]

%------------------------------------------------
\section{Bayesian Logistic Regression}

Model:
\[
P(Y_i=1\mid \beta)=\frac{e^{\beta x_i}}{1+e^{\beta x_i}},
\qquad
\beta\sim N(0,\tau^2)
\]

\begin{formulabox}
Posterior kernel:
\[
p(\beta\mid y,x)\propto
\left[
\prod_{i=1}^n
\frac{e^{y_i\beta x_i}}{1+e^{\beta x_i}}
\right]
\exp\left(-\frac{\beta^2}{2\tau^2}\right)
\]
\end{formulabox}

Odds ratio for one-unit increase in \(x\):
\[
e^\beta
\]

Approx.\ MAP = maximizer of log-posterior.

Exponentiate posterior quantiles of \(\beta\) to get credible interval for \(e^\beta\).

%------------------------------------------------
\section{Metropolis--Hastings}

Target density \(\pi(\theta)\), proposal \(q(\theta^\star\mid \theta)\).

\begin{formulabox}
Acceptance probability:
\[
\alpha(\theta,\theta^\star)=
\min\left\{
1,\,
\frac{\pi(\theta^\star)\,q(\theta\mid\theta^\star)}
{\pi(\theta)\,q(\theta^\star\mid\theta)}
\right\}
\]
\end{formulabox}

If proposal symmetric (\(q(\theta^\star\mid\theta)=q(\theta\mid\theta^\star)\)):
\[
\alpha(\theta,\theta^\star)=
\min\left\{1,\frac{\pi(\theta^\star)}{\pi(\theta)}\right\}
\]

Random-walk proposal:
\[
\theta^\star\sim N(\theta,s^2)
\]

Rule of thumb: moderate acceptance rate is usually good.

%------------------------------------------------
\section{Gaussian Mixture Model}

Latent variable form:
\[
z_i\mid \pi \sim Bernoulli(\pi),
\qquad
y_i\mid z_i,\mu_0,\mu_1,\sigma^2 \sim N(\mu_{z_i},\sigma^2)
\]

Priors:
\[
\pi\sim Beta(a,b),\quad
\mu_k\sim N(m_k,s_k^2),\quad
\sigma^2\sim Inv\text{-}Gamma(\alpha,\beta)
\]

\subsection{Full Conditional for \(z_i\)}

\begin{formulabox}
\[
P(z_i=1\mid \cdot)=
\frac{\pi\,\phi(y_i;\mu_1,\sigma^2)}
{(1-\pi)\phi(y_i;\mu_0,\sigma^2)+\pi\phi(y_i;\mu_1,\sigma^2)}
\]
\end{formulabox}

\subsection{Full Conditional for \(\pi\)}

If \(n_1=\sum_i z_i\), \(n_0=n-n_1\),
\[
\pi\mid z \sim Beta(a+n_1,b+n_0)
\]

\subsection{Full Conditional for \(\mu_k\)}

For \(k\in\{0,1\}\), let \(n_k=|\{i:z_i=k\}|\), \(\bar y_k\) sample mean in cluster \(k\).

\begin{formulabox}
\[
(\,s_k^\star\,)^{-2}=s_k^{-2}+\frac{n_k}{\sigma^2}
\]
\[
m_k^\star=(s_k^\star)^2
\left(
\frac{m_k}{s_k^2}+\frac{n_k\bar y_k}{\sigma^2}
\right)
\]
\[
\mu_k\mid \cdot \sim N(m_k^\star,(s_k^\star)^2)
\]
\end{formulabox}

\subsection{Full Conditional for \(\sigma^2\)}

Let
\[
SSE=\sum_{i=1}^n (y_i-\mu_{z_i})^2
\]

\begin{formulabox}
\[
\sigma^2\mid \cdot \sim Inv\text{-}Gamma\left(\alpha+\frac n2,\; \beta+\frac{SSE}{2}\right)
\]
\end{formulabox}

\subsection{One Gibbs Sweep}

Update in order:
\begin{enumerate}[leftmargin=14pt,itemsep=1pt]
\item all \(z_i\)
\item \(\pi\)
\item \(\mu_0\)
\item \(\mu_1\)
\item \(\sigma^2\)
\end{enumerate}

%------------------------------------------------
\section{Fast Problem-Solving Checklist}

\begin{theorembox}
\textbf{Common moves on exams}
\begin{itemize}[leftmargin=12pt,itemsep=1pt]
\item Stationary dist.: solve \(\pi P=\pi\) or use detailed balance.
\item Long-run fraction of time in state \(i\): use \(\pi(i)\).
\item Mean return time: use \(1/\pi(i)\).
\item Hitting probs.: write recursion + boundary conditions.
\item Expected hitting time: write \(m_i=1+\sum_j P_{ij}m_j\).
\item Absorbing chain: compute \(F=(I-Q)^{-1}\), then \(FR\).
\item Poisson counts: integrate \(\lambda(t)\).
\item Brownian drift MLE: net displacement / total time.
\item MH: posterior ratio \(\times\) proposal ratio.
\item Gibbs: identify full conditional families.
\end{itemize}
\end{theorembox}

\end{multicols}

\end{document}
